\section{Vietnamese-Specific OCR Challenges and Solutions}
Vietnamese presents a unique set of challenges for optical character recognition that extend well beyond those encountered in most Latin-script languages, and that are not adequately addressed by OCR systems designed for general multilingual use. While the deep learning architectures surveyed in the preceding section have demonstrated strong performance across a broad range of scripts and document types, their direct application to Vietnamese text without domain-specific adaptation consistently yields suboptimal results. These shortcomings stem from two interconnected factors that distinguish Vietnamese OCR from the general case.

The first factor is linguistic and orthographic in nature. Vietnamese employs a diacritical system of exceptional complexity, in which multiple combining marks can co-occur on a single character to encode both vowel quality and lexical tone simultaneously. This orthographic density produces a character inventory significantly larger than that of unaccented Latin-script languages, and introduces visual ambiguities that challenge both detection and recognition components of an OCR pipeline. The second factor is resource-related. Despite growing interest in Vietnamese document digitization, the landscape of publicly available annotated datasets and standardized benchmarks for Vietnamese OCR remains considerably less developed than the corresponding ecosystems for English, Chinese, or Arabic, constraining the training and rigorous evaluation of recognition models.

These two factors are closely interrelated: the orthographic complexity of Vietnamese makes the annotation process more demanding and error-prone, which in turn limits the scale and coverage of available datasets, which in turn impedes the development of models robust enough to handle the full range of Vietnamese diacritical variation. This section examines both dimensions in turn, first analyzing the specific diacritical and tonal characteristics of Vietnamese that create difficulties for OCR models, and then surveying the existing datasets and benchmarks that have been developed to support Vietnamese text recognition research, identifying the gaps that motivate the data preparation strategy adopted in this thesis.

\subsection{Diacritical Marks and Tonal Complexity}
Vietnamese presents a distinctively challenging target for optical character recognition due to its orthographic system, which overlays multiple layers of diacritical modification onto a Latin-derived base script. Unlike most Latin-script languages, where diacritics serve a limited and relatively predictable role, Vietnamese employs diacritics at two independent levels simultaneously: vowel quality markers that alter the phonetic identity of the base character, and tonal markers that specify one of six lexically contrastive tones \cite{Nguyen2012}. These two layers of modification can co-occur on a single character, producing composite glyphs that are visually dense and spatially compact, creating significant difficulties for both text detection and recognition systems.

At the level of vowel modification, Vietnamese augments several base vowels with diacritic strokes that fundamentally change their phonetic identity. The vowel \textit{a}, for instance, yields three distinct characters: \textit{a}, \textit{a with circumflex} (â), and \textit{a with breve} (ă), each representing a different phoneme and therefore a different lexical unit. Similar modification patterns apply to the vowels \textit{e}, \textit{o}, and \textit{u}, yielding a total of seventeen distinct vowel surface forms from seven base vowels. When tonal markers are superimposed on these modified vowels, the resulting character set expands substantially. Vietnamese defines six tones encoded through five diacritical marks, namely the grave, acute, hook above, tilde, and dot below, plus the absence of any tone mark for the level tone. Combining all vowel modifications with all tonal possibilities yields a Vietnamese character inventory that is several times larger than that of unaccented Latin-script languages, posing a direct challenge to recognition models trained on standard Latin benchmarks \cite{Nguyen2021}.

From the perspective of OCR model design, this orthographic complexity manifests in two principal failure modes. The first is character-level confusion, where visually similar composite glyphs are misclassified as one another due to the small spatial extent of distinguishing diacritical strokes relative to the base character body. At low image resolutions or under degraded scan conditions, the difference between a dot below and a hook above, or between a circumflex and a breve, can be reduced to a handful of pixels, well below the discriminative capacity of feature extractors trained on higher-resolution data \cite{Nguyen2021}. The second failure mode is segmentation ambiguity, where the spatial extent of a composite Vietnamese character overlaps with neighboring characters or with the ascender and descender zones of adjacent text lines, causing segmentation-based detectors to either merge distinct characters or fragment a single character across multiple segments \cite{Nguyen2021}. This segmentation problem is particularly acute in printed documents with tight inter-line spacing and diverse font styles, which fall squarely within the scope of this thesis.

Existing approaches to mitigating these challenges fall into two broad categories. The first encompasses character set redesign strategies, in which the model output vocabulary is explicitly constructed to include all valid Vietnamese composite characters as atomic units, rather than decomposing them into base characters and combining diacritics \cite{Nguyen2021}. This atomic treatment forces the recognition model to learn a direct mapping from visual glyph to full Unicode character, avoiding the error accumulation that arises when base character and diacritic predictions are made independently. The second category leverages linguistic constraints through language model integration, exploiting the fact that not all combinations of base character and diacritical marks are phonologically valid in Vietnamese, and that valid tone-vowel combinations follow predictable phonotactic rules that can be encoded as decoding constraints \cite{Nguyen2021}. The methodology adopted in this thesis draws on both categories, as detailed in Chapter 3.

\subsection{Existing Vietnamese OCR Datasets and Benchmark}

The availability of high-quality annotated datasets is a prerequisite for training and evaluating robust Vietnamese OCR systems, yet the Vietnamese OCR research community has historically been constrained by a limited and fragmented landscape of publicly available corpora. Compared to the rich benchmark ecosystems surrounding English and Chinese OCR, Vietnamese-specific datasets remain sparse in terms of both scale and diversity, a gap that has directly impeded the development of generalizable recognition models \cite{Nguyen2021}.

Among the earliest publicly available resources is the VOHTR2018 dataset \cite{Nguyen2018}, which targets offline handwritten Vietnamese text and provides word-level and line-level images annotated with ground truth transcriptions. While valuable for handwriting recognition research, VOHTR2018 does not address the printed document domain and its annotation granularity is insufficient for training text detection models that require bounding box supervision at the word or character level. Its annotations also do not generalize readily to scene text conditions, limiting its utility for the broader Vietnamese OCR task.

A more recent and comprehensive contribution is VinText \cite{Nguyen2021}, a Vietnamese scene text dataset comprising over 2,000 real images collected from diverse real-world environments including street signs, storefronts, menus, and printed documents. VinText provides word-level polygon annotations along with full transcriptions and constitutes the most linguistically diverse publicly available Vietnamese OCR benchmark to date. Its inclusion of natural scene images with complex backgrounds, arbitrary text orientations, and mixed font styles makes it substantially more challenging than synthetic corpora and more representative of deployment conditions. A key contribution of VinText is its dictionary-guided recognition framework, which explicitly incorporates Vietnamese lexical constraints to resolve diacritical ambiguities at inference time \cite{Nguyen2021}. However, its scale remains modest relative to large-scale multilingual benchmarks such as MLT \cite{Nayef2017} and LSVT \cite{Sun2019}, and its coverage of formal administrative document layouts is limited.

Across these datasets, several consistent gaps are apparent. First, no existing Vietnamese OCR benchmark provides sufficient coverage of the full range of document types encountered in administrative digitization workflows, including low-resolution scans, faded ink, and degraded print quality. Second, character-level annotation remains rare, with most datasets providing only word-level or line-level ground truth, which limits the ability to train and evaluate detection components at finer granularity. Third, the absence of a standardized evaluation protocol across datasets makes cross-paper comparison of reported results unreliable. These gaps collectively motivate the data preparation strategy adopted in this work, which combines synthetic generation with domain-specific augmentation to compensate for the scarcity of real annotated Vietnamese document images, as described in Chapter 3.
